\textit{Theorem 3 (conditional protection relative to an executable route).}
Let $\mathcal P$ and $\mathcal U_V$ be finite nonempty sets, let all supplied
route scores be finite, and let $\mu,\tau\ge0$.  Under
\eqref{eq:reference-difference} and \eqref{eq:protected-route}, either
$P_g=P_b$, or
\begin{equation}
\Delta(P_g,M)<-\mu-\tau
\quad\text{for every }M\in\mathcal U_V.
\label{eq:protected-score-bound}
\end{equation}
Fix a common realized environment and the same deterministic execution
rule for all routes.  Let $Q(P)$ be its cost under route $P$.  If some
$M_0\in\mathcal U_V$ and $\varepsilon_D\ge0$ satisfy
\begin{equation}
\left|Q(P)-Q(P_b)-\Delta(P,M_0)\right|\le\varepsilon_D
\quad\text{for all }P\in\mathcal P,
\label{eq:differential-error-condition}
\end{equation}
then
\begin{equation}
Q(P_g)-Q(P_b)\le\max\{0,\varepsilon_D-\mu-\tau\}.
\label{eq:protected-execution-bound}
\end{equation}
In particular, $\varepsilon_D\le\mu+\tau$ is sufficient for noninferiority
to $P_b$ in that environment.  Uniform absolute error
$|Q(P)-\widehat Q_t(P,M_0)|\le\varepsilon_Q$ implies
\eqref{eq:differential-error-condition} with
$\varepsilon_D=2\varepsilon_Q$.

\textit{Proof.}  The finite nonempty portfolio and finite scores give a
forecast matched minimizer $P_b$.  The same properties give each maximum
$U(P)$ and a minimizer $P_u$.  Because $\Delta(P_b,M)=0$ for every matrix,
$U(P_b)=0$ and $U(P_u)\le0$.  If the strict acceptance condition fails,
\eqref{eq:protected-route} returns $P_b$, whose score difference and
realized cost difference from itself are both zero.  If the condition
holds, then for every $M\in\mathcal U_V$,
\[
\Delta(P_u,M)\le U(P_u)<-\mu-\tau.
\]
This proves \eqref{eq:protected-score-bound}.  Applying
\eqref{eq:differential-error-condition} at $P_u$ and $M_0$ gives
\begin{align}
Q(P_u)-Q(P_b)
&\le\Delta(P_u,M_0)+\varepsilon_D \nonumber\\
&\le U(P_u)+\varepsilon_D
<-\mu-\tau+\varepsilon_D.
\end{align}
Combining this strict inequality with the zero difference in the fallback
branch proves \eqref{eq:protected-execution-bound}.  Under the absolute
error condition, define $e(P)=Q(P)-\widehat Q_t(P,M_0)$.  Then
\[
|Q(P)-Q(P_b)-\Delta(P,M_0)|
=|e(P)-e(P_b)|\le2\varepsilon_Q,
\]
which proves the final assertion.  \hfill $\square$

The same matrix must be used in both terms of
\eqref{eq:reference-difference}.  For example, suppose the reference
scores at two matrices are $(100,0)$ and another route has scores
$(99,10)$.  Subtracting their separate maxima gives $-1$, although the
second matrix gives a cost increase of $10$.  The paired maximum is
$10$, so \eqref{eq:protected-route} does not accept that route.

Theorem 3 is an algebraic statement about the supplied finite scores.
For computed differences $\widetilde\Delta$, suppose additionally that
$|\widetilde\Delta(P,M)-\Delta(P,M)|\le\varepsilon_{\rm num}$ uniformly.
A route accepted using the computed maximum then satisfies
$\Delta(P_g,M)<-\mu-\tau+\varepsilon_{\rm num}$ for each matrix, by adding
the difference error to the strict acceptance inequality.  Repeating the
proof gives \eqref{eq:protected-execution-bound} with
$\varepsilon_D+\varepsilon_{\rm num}$ in place of $\varepsilon_D$.
If each computed score has absolute error at most $\xi$ and the subtraction
has additional error at most $\chi$, the triangle inequality gives
$\varepsilon_{\rm num}\le2\xi+\chi$.  The chosen $\tau$ alone does not
establish either error bound.

Returning $P_b$ means executing its route rather than merely assigning its
policy label.  Given identical initial states and future inputs, the two
executions have identical service allocations at the first step.  Their
packet and electrical operations use the same inputs and deterministic
acceptance rule, and their completed jobs are identical.  The service and
state recursions therefore give identical next states.  Induction repeats
this argument through the finite horizon, proving equality of the returned
reference execution and its cost.  Numerical comparisons additionally check
the realized records.  Neither finite portfolios nor those checks imply a
small uniform $\varepsilon_D$ on new environments.  Reference selection and
its acceptance branch remain outside the continuous map of Corollary 1.
